\documentclass[11pt]{article}
\usepackage[margin=1.05in]{geometry}
\usepackage{amsmath,amssymb,booktabs,xcolor,hyperref,microtype}
\usepackage[T1]{fontenc}
\usepackage{lmodern}
\hypersetup{colorlinks,linkcolor=blue!50!black,citecolor=blue!50!black,urlcolor=blue!50!black}
\newcommand{\R}{\mathbb{R}}
\newcommand{\Om}{\Omega}
\newcommand{\Gam}{\Gamma}
\title{L-sweeps for the sparsified Lippmann--Schwinger equation}
\date{Notes on the \texttt{Lsweeps} branch of \texttt{SparsifyAndSweep.jl}, September 2026}
\author{}
\begin{document}
\maketitle
\vspace{-2em}

\section{Setting, and what limits the layered sweep}

The package solves the two-dimensional Lippmann--Schwinger (LS) equation for the
Helmholtz scattering problem,
\begin{equation}
  u(x) + \omega^2\!\int_\Om G(x-y)\,m(y)\,u(y)\,dy \;=\; -\omega^2\!\int_\Om G(x-y)\,m(y)\,u_I(y)\,dy,
  \qquad G(x)=\tfrac{i}{4}H^{(1)}_0(\omega|x|),
  \label{eq:ls}
\end{equation}
with $m=1-1/c^2$ compactly supported in $\Om=(0,1)^2$, discretised by a Nystr\"om rule on a
uniform grid with $h=\lambda/8$ and a corrected diagonal weight. The dense system
$(I+\omega^2 KM)u=g$ is applied by a circulant embedding and solved by GMRES; the
preconditioner is that of Liu and Ying~\cite{LiuYing2018}, in two stages.

\emph{Sparsify.} A single $9$-point stencil $\alpha$ is fitted so that $\alpha^*$
annihilates the columns of $K$ centred outside the $3\times3$ neighbourhood $\mu$,
$\min_{\|\alpha\|=1}\|\alpha^*K_{\mu,\mu^c}\|$, and $\beta^*=\alpha^*K_{\mu\mu}$. Applying
$\alpha^*$ to each row of the dense system gives a sparse surrogate $H\tilde u=f$, with rows
$\alpha^*\tilde u_{\mu_i}+\omega^2\beta^*[m\tilde u]_{\mu_i}$ on the grid and rows
$\gamma_i^*\tilde u_{\mu_i}$ on $b$ layers of perfectly matched layer (PML) around it. The
$\gamma$ stencils are fitted to annihilate complex-stretched plane waves
$e^{i\kappa\,r\cdot x^\sigma}$ at the local frequency $\kappa=\omega\sqrt{1-m(x)}$, so they
are valid \emph{inside} the medium, not only on the outer boundary. Inverting $H$ exactly
costs $3$ GMRES iterations at every frequency tried; everything above $3$ comes from the
second stage.

\emph{Sweep.} $H$ is block tridiagonal over slices of width $b$ along $x_1$. The moving-PML
sweep approximates each Schur complement by a local problem on the slice padded with
$\gamma$ rows, and applies the resulting block $LDU$ factorisation from both ends towards
the middle. This is the bottleneck for parallelism: the application is a linear
recurrence over the $\ell=n/b$ slices. Two fronts give concurrency two, and no amount of
hardware shortens the recurrence. Measured on this code, the intrinsic speed-up of one
application over slices is $1.8$--$1.95$ in 2D and $1.6$ in 3D, and for a single
right-hand side the whole solve cannot gain more than about $4\times$ from slice
parallelism, because setup is half the time and the application is capped at two.

\section{L-sweeps in a nutshell}

Taus, Zepeda-N\'u\~nez, Hewett and Demanet~\cite{Taus2020} remove the recurrence by
changing the decomposition. The domain is cut into a Cartesian array of $q\times r$ boxes
$\Om_{ij}$, each extended by a thin layer and a PML into a local problem with matrix
$A_{ij}$. Information moves between boxes through \emph{polarized wavefields}: given the
values of a field on the two grid layers $\Gam_1\,|\,\Gam_2$ on either side of an
interface, one local solve reproduces the field beyond the interface and vanishes before
it. Sources anywhere are handled in three stages: a local solve in every box; four
\emph{cardinal} sweeps, right, left, up and down, along the rows and columns of boxes;
and four \emph{diagonal} sweeps, corner to corner, in which the boxes on one
anti-diagonal are updated together from an L-shaped trace. Each sweep is a chain, but the
chains are independent: the $q$ rows in a horizontal sweep, the $r$ columns in a vertical
one, the $\min(q,r)$ boxes of an anti-diagonal in a diagonal one. With one rank per row
of boxes and boxes of fixed size in wavelengths, the application costs $O(N/p)$ on
$p=O(n)$ ranks, with only nearest-neighbour messages of one interface layer each. For a
constant medium and a thick enough PML the preconditioner is a direct solver (one
iteration); in general the iteration count grows like $\log\omega$.

The price is work: nine local solves per box per application against two slab solves
per slice for the layered sweep, and up to four shifted decompositions with a partition
of unity when a source sits on the skeleton of the decomposition. What is bought is a
critical path of $O(q)$ box solves of $O(1)$ size instead of $O(\ell)$ slab solves of
$O(n)$ size.

\section{The transfer to the sparsified system}

\subsection{The identity is algebraic, and $H$ satisfies its hypothesis}

Order the unknowns of a local matrix $A$ as $(\Om_1,\Gam_1,\Gam_2,\Om_2)$, where $\Gam_1$
and $\Gam_2$ are the layers adjacent to the interface on each side. If the stencil
couples only adjacent layers, $A$ is block tridiagonal in this ordering. Let $v$ satisfy
$Av=f$ with $f=0$ on $\Gam_2\cup\Om_2$. Then
\begin{equation}
  A\,U = \bigl(0,\;A_{\Gam_1\Gam_2}v_{\Gam_2},\;-A_{\Gam_2\Gam_1}v_{\Gam_1},\;0\bigr)
  \quad\Longrightarrow\quad U=(0,\,0,\,v_{\Gam_2},\,v_{\Om_2}),
  \label{eq:pol}
\end{equation}
which one checks row block by row block: the $\Om_1$ and $\Gam_1$ rows hold because $U$
vanishes on $\Om_1\cup\Gam_1$, the $\Gam_2$ row is the equation $(Av)_{\Gam_2}=0$
rearranged, and the $\Om_2$ row is $(Av)_{\Om_2}=0$. Nothing in this uses symmetry, a
Green's identity, or the fact that $A$ discretises a differential operator. The sparsified
$H$ has a compact $3\times3$ stencil, so one layer on each side suffices, exactly as for the
five-point Laplacian of the paper. On the assembled $H$, with exact traces, the identity
holds to $10^{-13}$ for straight interfaces and for the L-shaped interface of a
quarter-space.

The precedent is worth naming: Zepeda-N\'u\~nez and Zhao~\cite{ZepedaZhao2016} already
combined Ying's sparsifying preconditioner with the layered method of polarized traces for
the 2D LS equation. The obstacle they met is instructive. Since the sparsified matrix does
not come from a differential operator, no absorbing boundary condition was available for
its local problems; they built one from a smooth cut-off of $m$ with a complex shift, and
found that one-directional sweeps then needed $O(\sqrt\omega)$ iterations, which alternating
directions brought back to $O(\log\omega)$. The 2018 stencils $\gamma$, fitted at the local
frequency, are precisely the absorbing rows that were missing.

\subsection{Local box problems from the stencils alone}

A box $\Om_{ij}$ is extended by two unstretched layers and then by $p$ layers of PML on
every side that has a neighbour; boxes on the outer boundary keep the global PML there. A
row of the local matrix is decided by the stretching increments $\Delta\sigma_d$ of the
point on each axis: zero increments give the $\alpha/\beta$ row of $H$, and any non-zero
increment gives a $\gamma$ row fitted at $\kappa=\omega\sqrt{1-m}$ with those increments,
corner points included. The per-axis profile is the two-sided quadratic ramp already used
for the global PML, with $\pm\infty$ edges on sides without a pad. Consequently a local
matrix needs only $\alpha$, $\beta$, $m$ and the geometry, never $H$: every rank can build
and factorise exactly the boxes it owns, which is the embarrassingly parallel setup the
paper relies on.

Two conventions turned out to matter, each by three orders of magnitude, and both were
found by getting them wrong first.
\begin{itemize}
\item \textbf{Where the stretching starts.} The zero of the PML profile must sit on the
\emph{first} pad row, so that the stretching begins one row inside the pad; this is what
the moving PML of the layered sweep does. A closed $64^2$ box with a source inside then
reproduces the free-space solution to $4.4\cdot10^{-5}$ in a constant medium and
$6.6\cdot10^{-4}$ inside the lens, independently of the pad width and of the source
position. Placing the zero on the last interior row instead gives $14$--$15\%$.
\item \textbf{The corner of an L-shaped trace.} The corner degree of freedom belongs to
the diagonal neighbour and must take its value from that box. Taking it from a cardinal
neighbour's polarized field, which vanishes there by construction, leaves a
pad-independent error near $1\%$ in the diagonal box; with the corner from its owner the
diagonal step is as accurate as the source box.
\end{itemize}

\subsection{Sources everywhere without windows}

In the preconditioner the right-hand side is $f=Sr$ for a Krylov residual $r$, so it is
non-zero at every grid point. The paper's fourth scenario, with four shifted
decompositions and a partition of unity, exists because a continuous source can straddle
an interface. On a grid whose interfaces run \emph{between} grid lines no degree of freedom
lies on the skeleton, so the third scenario, arbitrary sources not intersecting the
skeleton, already covers every right-hand side and the factor four is not paid.

What must be paid instead is bookkeeping against double counting. The field of a box is
kept as nine components: its own local solution $w^{0}$, what arrived from each cardinal
direction $w^{\rightarrow},w^{\leftarrow},w^{\uparrow},w^{\downarrow}$, and what arrived
along each diagonal $w^{\nearrow},w^{\nwarrow},w^{\searrow},w^{\swarrow}$. The trace
handed on sums exactly the components that carry the sources of the relevant half-plane or
quadrant. For the upward step from $(i,j)$ into $(i{+}1,j)$ the trace is $w^0+w^\uparrow$
of $(i,j)$. For the north-east diagonal step into $(i,j)$, the horizontal arm carries
$w^\rightarrow+w^\nearrow$ of the box below, the vertical arm $w^\uparrow+w^\nearrow$ of
the box to the left, and the corner $w^0+w^\rightarrow+w^\uparrow+w^\nearrow$ of the
diagonal box. The final field is the sum of all nine components on the owned points.

\subsection{What does not transfer}

\emph{The one-iteration property moves inside.} In the paper the preconditioner inverts the
matrix being solved, so it becomes exact as the PML thickens. Here it inverts a
surrogate, and the outer iteration runs on the dense operator: the exact $H^{-1}$ leaves
$3$ iterations, and thickening the absorber of the layered sweep from one to four
wavelengths brings the count from $6$ down to $4$, not to $1$. The property is
available where the two-level scheme of~\cite{ZepedaZhao2016} needs it: a few GMRES steps
on $H$ inside each application, preconditioned by the sweep, bring the outer count to the
floor of $3$. Because such a preconditioner is not a fixed linear operator, the outer
solver must be a flexible GMRES that stores the preconditioned basis; with plain GMRES the
iteration reports convergence while the true residual stalls near $10^{-6}$.

\emph{The outer operator is global.} The paper's Krylov products are sparse and local.
Ours are dense: a global FFT with an all-to-all transpose, or a fast multipole product
with its own tree layout. In the present version every rank keeps the full vectors and
runs the outer iteration replicated, so the sweep is distributed but the dense product is
not; distributing it is the next step, and the row-based box ownership is already a
decomposition into horizontal bands that a slab FFT or a band-partitioned fast method can
share. With a fast multipole product the FFT disappears from the solver altogether, since
the preconditioner itself uses none.

\section{Implementation: one dataflow for threads, MPI and both}

The sweep is expressed as a dependency graph rather than as nested loops. A \emph{node}
is a pair (box, direction): the arrival of one component in one box, with direction
$0$ for the box's own solve. A node has \emph{contributions}, each a set of trace points
whose values are the sum of named components of one neighbouring box, available once the
nodes producing those components have fired. All of this is geometric and is computed on
every rank at setup. At solve time a node fires when its contributions are present; firing
it runs the polarized solve~\eqref{eq:pol} and releases the contributions that depend on
it. If source and target boxes are on the same rank the values are copied; otherwise
they travel as one message, tagged by target node and contribution, into a receive posted
at the start of the solve. Within a rank a pool of tasks executes ready nodes, so rows,
columns and anti-diagonals overlap as far as the dependencies allow; the calling task
polls the outstanding receives. Boxes are assigned to ranks row by row, the paper's
assignment, so horizontal sweeps never communicate and columns and diagonals pipeline
across ranks; an arbitrary ownership matrix is accepted. The owned parts of the result
are combined by one all-reduce.

Every node's arithmetic depends only on its inputs, summed in a fixed order, so
sequential, threaded and distributed runs agree bitwise. The tests check this for one
against eight worker tasks in one process, and for three ranks of two threads against one
process. The code is in \path{src/lsweep.jl}; the MPI backend is the weak-dependency
extension \path{ext/SparsifyAndSweepMPIExt.jl}, which serialises MPI calls behind a lock
so that any thread level from \emph{serialized} up suffices. It is selected with
\path{SweepPreconditioner(P; mode = :lsweep, nbox, pad, comm, inner)}.

\section{Numerical evidence}

\paragraph{One transmission.} A minimal $2\times2$ decomposition with $64$-point boxes, the
corner at the centre of the converging lens ($m\approx-1$ there), source in the lower-left
box; relative error of each box against the global solution of $H$:
\begin{center}\small
\begin{tabular}{llrrr}
\toprule
medium, $\omega/2\pi$ & pad & source box & straight steps & diagonal step \\
\midrule
constant, 32 & $1\lambda$ & $8\cdot10^{-5}$ & $7.5\cdot10^{-4}$ & $1.3\cdot10^{-4}$ \\
constant, 32 & $2\lambda$ & $8\cdot10^{-5}$ & $3.0\cdot10^{-4}$ & $1.3\cdot10^{-4}$ \\
lens, 32     & $1\lambda$ & $4\cdot10^{-4}$ & $3.8\cdot10^{-3}$ & $4.5\cdot10^{-4}$ \\
lens, 32     & $2\lambda$ & $2.6\cdot10^{-4}$ & $1.5\cdot10^{-3}$ & $1.7\cdot10^{-4}$ \\
lens, 64     & $1\lambda$ & $1.1\cdot10^{-3}$ & $1.8\cdot10^{-2}$ & $2.0\cdot10^{-3}$ \\
lens, 64     & $2\lambda$ & $4.7\cdot10^{-4}$ & $8.5\cdot10^{-3}$ & $4.6\cdot10^{-4}$ \\
lens, 64     & $3\lambda$ & $4.2\cdot10^{-4}$ & $4.1\cdot10^{-3}$ & $5.2\cdot10^{-4}$ \\
lens, 128    & $1\lambda$ & $1.0\cdot10^{-3}$ & $1.1\cdot10^{-2}$ & $1.5\cdot10^{-3}$ \\
lens, 128    & $2\lambda$ & $7.8\cdot10^{-4}$ & $5.2\cdot10^{-3}$ & $1.1\cdot10^{-3}$ \\
lens, 128    & $3\lambda$ & $7.8\cdot10^{-4}$ & $3.5\cdot10^{-3}$ & $1.1\cdot10^{-3}$ \\
\bottomrule
\end{tabular}
\end{center}
The straight steps halve with each doubling of the absorber, the signature of a working
PML; the floor near $10^{-4}$ is the accuracy of the fitted stencils themselves; the
diagonal step is as accurate as the source box.

\paragraph{One full application} against $H^{-1}f$ for a random $f$ supported on all of
$\Om$ ($n=127$, $\omega/2\pi=16$): in a constant medium between $1.8\cdot10^{-4}$ and
$6.3\cdot10^{-4}$ for $2\times2$ to $4\times4$ boxes, halving from one to two wavelengths
of pad; in the lens between $2\cdot10^{-3}$ and $3\cdot10^{-2}$ depending on where the
interfaces cut it, and nearly independent of the pad: as in \S4 of~\cite{Taus2020}, the local problems do not see
reflections between boxes, and the outer iteration supplies them.

\paragraph{As the preconditioner.} Converging lens, $h=\lambda/8$, global PML width
$b=8$, GMRES(20) to $10^{-6}$, one process with $16$ threads on a $16$-core desktop.
Boxes of about $64$ points, each padded by $8$ points --- one wavelength at this $h$ ---
unless the row says otherwise, and UMFPACK's own factors throughout; the compact factors
come later in this section. Each row was measured in a process of its own, so the peak
resident set belongs to that row alone:
\begin{center}\footnotesize\setlength{\tabcolsep}{5pt}
\begin{tabular}{rlrrcrr}
\toprule
$\omega/2\pi$ & inner solver & setup & apply & iterations & solve & peak \\
\midrule
 32 & layered sweep      & 1.08 s & 0.036 s & 5 & 0.25 s & 0.7 GB \\
    & L-sweeps           & 0.50 s & 0.024 s & 4 & 0.18 s & 0.9 GB \\
    & L-sweeps, inner 6  & 0.46 s & 0.075 s & 3 & 0.27 s & 0.9 GB \\
 64 & layered sweep      & 1.99 s & 0.131 s & 6 & 1.03 s & 1.4 GB \\
    & L-sweeps           & 1.24 s & 0.115 s & 5 & 0.68 s & 1.8 GB \\
    & L-sweeps, inner 6  & 1.25 s & 0.398 s & 3 & 1.09 s & 1.9 GB \\
128 & layered sweep      & 6.22 s & 0.551 s & 7 & 4.55 s & 4.1 GB \\
    & L-sweeps           & 3.33 s & 0.393 s & 7 & 3.61 s & 5.6 GB \\
    & L-sweeps, pad 16   & 5.03 s & 0.595 s & 6 & 4.26 s & 7.2 GB \\
    & L-sweeps, inner 6  & 3.30 s & 1.531 s & 3 & 4.95 s & 5.7 GB \\
256 & layered sweep      & 29.36 s & 2.059 s & 9 & 23.99 s & 15.1 GB \\
    & L-sweeps           & 11.10 s & 1.554 s & 14 & 28.34 s & 20.5 GB \\
    & L-sweeps, pad 16   & 15.36 s & 2.239 s & 10 & 27.95 s & 27.0 GB \\
    & L-sweeps, pad 24   & 22.53 s & 3.042 s & 8 & 29.70 s & 34.9 GB \\
    & L-sweeps, inner 6  & 11.14 s & 9.245 s & 3 & 27.98 s & 20.5 GB \\
\bottomrule
\end{tabular}
\end{center}
Despite nine solves per box against two per slice, L-sweeps sets up between $2.0$ and
$2.6$ times faster at every size --- at $\omega/2\pi=256$ it factorises $1024$ local
problems of $7140$ unknowns where the layered sweep factorises $256$ slabs of up to
$51\,625$ --- and on time to solution with setup included it wins throughout, though the
margin narrows steadily: $1.95\times$ at $\omega/2\pi=32$, then $1.57$, $1.55$ and
$1.35\times$ at $256$. The mechanism behind both facts is not the one the table suggests,
so it is worth measuring rather than assuming.

\paragraph{What threads, and what does not.} \path{src/sweep.jl} contains no threading
constructs at all. The layered sweep is serial in both of its stages: its application is a
sequential recurrence, which is inherent to the method, but its factorisation loop over
slices writes only into per-slice arrays, so it is embarrassingly parallel and simply is
not threaded. Speedups from $1$ to $16$ threads, with BLAS pinned to the Julia thread count
so that a one-thread run cannot quietly use all sixteen cores inside UMFPACK's dense
kernels:
\begin{center}\small
\begin{tabular}{rcccc}
\toprule
& \multicolumn{2}{c}{layered sweep} & \multicolumn{2}{c}{L-sweeps} \\
\cmidrule(lr){2-3}\cmidrule(lr){4-5}
$\omega/2\pi$ & factorisation & application & factorisation & application \\
\midrule
 32 & $1.03\times$ & $0.97\times$ & $1.10\times$ & $4.17\times$ \\
 64 & $1.03\times$ & $1.05\times$ & $2.00\times$ & $6.40\times$ \\
128 & $0.96\times$ & $1.01\times$ & $2.30\times$ & $6.38\times$ \\
256 & $0.93\times$ & $0.98\times$ & $3.06\times$ & $6.39\times$ \\
\bottomrule
\end{tabular}
\end{center}
The layered columns are flat, and slightly below one at $256$. The dense operator is serial
too, for a duller reason: the package never calls \path{FFTW.set_num_threads}, so the FFT
convolution applied once per GMRES iteration costs the same at every thread count.

The layered sweep is fast regardless because it does far less work. On one core at
$\omega/2\pi=256$ one layered application costs $1.97$ s against $9.01$ s for one L-sweeps
application, a factor of $4.6$; taking each at its own iteration count, the layered sweep
reaches the answer in $22.6+9\cdot1.97=40.3$~s against $17.9+14\cdot9.01=144.1$~s, which
is $3.6$ times faster. Serial factorisation costs are close, $22.6$~s against $17.9$~s. So
every L-sweeps advantage in the preceding table is bought with the other fifteen cores: the
comparison is between a parallel implementation and a serial one, and what L-sweeps really
trades is about $4.6\times$ more application work for parallelism the layered
implementation does not yet exploit.

\paragraph{The absorber has to grow with the frequency.} What narrows the margin is the
iteration count, at a pad held fixed in wavelengths:
\begin{center}\small
\begin{tabular}{rcccc}
\toprule
$\omega/2\pi$ & layered sweep & pad $1\lambda$ & pad $2\lambda$ & pad $3\lambda$ \\
\midrule
 32 & 5 & 4 & 4 & 4 \\
 64 & 6 & 5 & 5 & 5 \\
128 & 7 & 7 & 6 & 6 \\
256 & 9 & 14 & 10 & 8 \\
\bottomrule
\end{tabular}
\end{center}
Some growth is expected: the paper reports $N_{\mathrm{it}}=O(\log\omega)$, which over
these four frequencies would read roughly $4,5,6,7$. The $3\lambda$ column, $4,5,6,8$,
about recovers that; the $1\lambda$ column outruns it. The natural explanation
--- that the reflection each absorber leaves accumulates over the $q\propto\omega$
interfaces a wave crosses --- is wrong, and was tested: at $\omega/2\pi=256$ and a
$1\lambda$ pad, boxes of $64$, $96$ and $128$ points ($q=32$, $22$, $16$) need $14$, $14$
and $15$ iterations, and at $2\lambda$ the $64$- and $128$-point boxes both need $10$.
Halving the interfaces changes nothing. And $128$-point boxes at $256$ against $64$-point
boxes at $128$ --- the same $q=16$, the same pad --- need $15$ against $7$. What sets the
count is the pad in wavelengths and the frequency itself, the size of $\Om$ in
wavelengths; why propagation distance should matter at fixed per-interface absorption is
open. Widening the absorber to $3\lambda$ at
$\omega/2\pi=256$ restores the count to $8$, below the layered sweep, at the price of a
local problem of $13\,572$ unknowns instead of $7140$: setup doubles, the application doubles,
and the peak grows to $34.9$ GB. For one right-hand side the narrow pad remains the better
trade. For many right-hand sides, where the setup amortises and the solve phase alone
decides, the layered sweep overtakes every L-sweeps variant with UMFPACK factors at $256$;
the compact factors of the next paragraph reverse that.

\paragraph{Half the bytes.} If the application is bound by the bytes it streams, the
lever is the factors themselves. Each box's UMFPACK factorisation is now re-stored, once
at setup, as $L$ and $U$ in single precision with 32-bit indices --- $12$ bytes per
nonzero against UMFPACK's $24$ --- together with the row scaling and permutations that
reproduce its solve, $(R_sA)[p,q]=LU$. The triangular solves run in Julia, accumulate in
double precision, allocate nothing, and skip the zeros a trace right-hand side leaves in
the running vector. The per-box solve agrees with UMFPACK's to $10^{-6}$, the
preconditioner is unchanged to the digits shown, and the result still agrees bitwise
across worker counts. At $\omega/2\pi=256$:
\begin{center}\footnotesize\setlength{\tabcolsep}{5pt}
\begin{tabular}{lrrcrr}
\toprule
$64$-point boxes unless stated & setup & apply & iterations & solve & peak \\
\midrule
pad 8, UMFPACK factors           & 10.27 s & 1.563 s & 14 & 28.61 s & 20.4 GB \\
pad 8, compact                   & 12.13 s & 1.089 s & 14 & 19.44 s & 14.5 GB \\
pad 16, compact                  & 16.82 s & 1.472 s & 10 & 18.36 s & 18.2 GB \\
pad 16, 128-point boxes, compact & 17.93 s & 1.069 s & 10 & 14.85 s & 15.8 GB \\
layered sweep                    & 29.36 s & 2.059 s &  9 & 23.99 s & 15.1 GB \\
\bottomrule
\end{tabular}
\end{center}
The application is $1.44\times$ faster and the peak $29\%$ lower, for two seconds of
extraction at setup; the gain is short of the $2\times$ in bytes because the solves are
latency-bound per core as well as bandwidth-bound in aggregate. The solve phase is again
ahead of the layered sweep's, $19.4$~s against $24.0$~s, and time to solution is $31.6$~s
against $53.4$. Larger boxes bought no iterations but do buy bytes --- the padded box's
redundancy falls from $1.72\times$ at $64$ points to $1.34\times$ at $128$ --- which is
why $128$-point boxes with a $2\lambda$ pad give the fastest solve phase measured,
$14.9$~s. Two smaller changes to the setup went in alongside. BLAS is pinned to one
thread while the boxes factorise, since its threads only spun at their barriers ($92$
CPU-seconds of $150$); the pin is unconditional, because the BLAS thread count changes the
factors in their last bits and one-worker and sixteen-worker runs are meant to agree
bitwise, and it is the same fix as the \path{OPENBLAS_NUM_THREADS=1} of the next table,
applied inside. And the local matrix is assembled straight into CSC rather than through a
coordinate list and \path{sparse()}, bitwise the same matrix for two thirds of the
allocation. Together they take the setup from $11.1$ to $10.3$~s here.

\paragraph{Memory.} A box padded by one wavelength carries $7140$ unknowns against the
$4160$ it owns, a $1.7\times$ redundancy. With UMFPACK's factors the measured peak is
about $1.35\times$ the layered sweep's at every size; with the compact ones it is $14.5$
GB against $15.1$ at $\omega/2\pi=256$, below it. At $\omega/2\pi=512$
($N=1.7\cdot10^{7}$) neither method fits the $61$ GB of this machine: both were stopped by
a watchdog at $40$ GB while still in setup, and scaling the $256$ peaks by the fourfold
growth in $N$ puts them near $65$ GB for the layered sweep and $58$ for L-sweeps with
compact factors --- closer, but not this machine. Distributing the boxes is what lifts
that ceiling, which is the subject of the next table.

\paragraph{Threads, ranks and both.} Setup and one application at $\omega/2\pi=256$
($1024$ boxes), best of three, ranks bound to disjoint cores, timed as the slowest rank:
\begin{center}\small
\begin{tabular}{lrrr}
\toprule
configuration & setup & apply & peak per rank \\
\midrule
1 process $\times$ 16 threads & 6.80 s & 1.43 s & 15.1 GB \\
4 ranks $\times$ 4 threads    & 4.54 s & 1.62 s & 4.7 GB \\
8 ranks $\times$ 2 threads    & 4.77 s & 1.98 s & 2.9 GB \\
16 ranks $\times$ 1 thread    & 4.48 s & 1.69 s & 2.1 GB \\
\bottomrule
\end{tabular}
\end{center}
The messages are two grid layers per interface and hide behind the local solves: the
application is flat across every configuration. Setup is flat too across the distributed
ones and half again slower with threads alone, which is the sparse LU's poor threading
showing through once more. The column that matters is the last: memory per rank falls
sevenfold from one process to sixteen, because the boxes are partitioned and only the
grid-sized vectors and the kernel are replicated. That is what puts problems larger than
one node can hold --- $\omega/2\pi=512$ among them --- within reach of a cluster.

Two configuration traps are worth recording, both of which cost a factor of three. Without
core binding the $4\times4$ run applies in $1.18$ s, because the launcher pins each rank to
a single core and its threads then share it. And with OpenBLAS left at its default thread
count, the $4\times4$ setup takes $13.5$ s instead of $4.5$ s: every rank starts a full
thread pool inside its four cores and the pools oversubscribe. The numbers above are with
\path{OPENBLAS_NUM_THREADS=1}, which is the right setting whenever more than one rank
shares a node.

\paragraph{A constant-contrast disc, against the Mie series.} A disc of radius $1/4$ at
$(1/2,1/2)$ --- the unit disc after rescaling lengths by four, so $\omega R=\omega/4$ ---
with velocity $0.8$ inside (index $1.25$) and $1$ outside, under the same plane wave and
otherwise the setup of the table above. The package solves for the scattered field, and
the Mie series gives it exactly; every row at a frequency solves the same discrete system
to $10^{-6}$, so the Mie error is the discretisation's and is the same on every row.
Below $\omega/2\pi=16$ the $64$-point box is the whole grid and the L-sweeps rows are the
exact solve of $H$ in single precision; ``exact solve of $H$'' is its UMFPACK LU.
\begin{center}\scriptsize\setlength{\tabcolsep}{2.2pt}
\begin{tabular}{rlrrrrrrrrrrr}
\toprule
$\omega/2\pi$ & inner solver & setup & apply & outer & sweeps & fwd & Hmv & $t_{\rm sw}$ & $t_{\rm fwd}$ & $t_{\rm Hmv}$ & solve & peak \\
\midrule
2 & layered sweep & 0.33 s & 0.000 s & 4 & 5 & 4 & 0 & 0.0000 s & 0.0000 s & 0.0000 s & 0.02 s & 0.5 GB \\
 & one box (= exact $H$) & 0.09 s & 0.000 s & 2 & 3 & 2 & 0 & 0.0000 s & 0.0000 s & 0.0000 s & 0.02 s & 0.6 GB \\
 & one box, inner 6 & 0.10 s & 0.001 s & 2 & 4 & 3 & 2 & 0.0005 s & 0.0000 s & 0.0000 s & 0.00 s & 0.6 GB \\
4 & layered sweep & 0.49 s & 0.001 s & 6 & 7 & 6 & 0 & 0.0010 s & 0.0000 s & 0.0000 s & 0.02 s & 0.5 GB \\
 & one box (= exact $H$) & 0.11 s & 0.000 s & 3 & 4 & 3 & 0 & 0.0000 s & 0.0000 s & 0.0000 s & 0.02 s & 0.6 GB \\
 & one box, inner 6 & 0.10 s & 0.001 s & 3 & 6 & 4 & 3 & 0.0005 s & 0.0000 s & 0.0000 s & 0.00 s & 0.6 GB \\
8 & layered sweep & 0.63 s & 0.003 s & 9 & 10 & 9 & 0 & 0.0030 s & 0.0001 s & 0.0000 s & 0.04 s & 0.6 GB \\
 & one box (= exact $H$) & 0.13 s & 0.001 s & 3 & 4 & 3 & 0 & 0.0010 s & 0.0001 s & 0.0000 s & 0.02 s & 0.6 GB \\
 & one box, inner 6 & 0.14 s & 0.003 s & 3 & 6 & 4 & 3 & 0.0015 s & 0.0001 s & 0.0000 s & 0.02 s & 0.6 GB \\
16 & layered sweep & 0.75 s & 0.009 s & 12 & 13 & 12 & 0 & 0.0090 s & 0.0005 s & 0.0001 s & 0.14 s & 0.6 GB \\
 & L-sweeps pad 8 & 0.22 s & 0.003 s & 12 & 13 & 12 & 0 & 0.0030 s & 0.0005 s & 0.0001 s & 0.09 s & 0.7 GB \\
 & L-sweeps pad 8, inner 6 & 0.22 s & 0.058 s & 3 & 24 & 4 & 21 & 0.0073 s & 0.0005 s & 0.0001 s & 0.08 s & 0.7 GB \\
32 & exact solve of $H$ & 0.47 s & 0.018 s & 3 & 4 & 3 & 0 & 0.018 s & 0.0021 s & 0.0005 s & 0.11 s & 0.7 GB \\
 & layered sweep & 1.05 s & 0.039 s & 16 & 17 & 16 & 0 & 0.039 s & 0.0021 s & 0.0005 s & 0.62 s & 0.8 GB \\
 & L-sweeps pad 8 & 0.47 s & 0.011 s & 19 & 20 & 19 & 0 & 0.011 s & 0.0021 s & 0.0005 s & 0.40 s & 1.2 GB \\
 & L-sweeps pad 8, inner 6 & 0.48 s & 0.127 s & 5 & 40 & 6 & 35 & 0.016 s & 0.0021 s & 0.0005 s & 0.54 s & 1.1 GB \\
64 & exact solve of $H$ & 1.60 s & 0.069 s & 4 & 5 & 4 & 0 & 0.069 s & 0.0092 s & 0.0029 s & 0.47 s & 1.4 GB \\
 & layered sweep & 1.98 s & 0.136 s & 34 & 36 & 35 & 0 & 0.136 s & 0.0092 s & 0.0029 s & 5.38 s & 1.5 GB \\
 & L-sweeps pad 8 & 1.23 s & 0.066 s & 32 & 34 & 33 & 0 & 0.066 s & 0.0092 s & 0.0029 s & 2.72 s & 1.7 GB \\
 & L-sweeps pad 8, inner 6 & 1.20 s & 0.481 s & 6 & 48 & 7 & 42 & 0.060 s & 0.0092 s & 0.0029 s & 3.13 s & 1.9 GB \\
128 & exact solve of $H$ & 7.06 s & 0.304 s & 5 & 6 & 5 & 0 & 0.304 s & 0.055 s & 0.012 s & 2.19 s & 4.4 GB \\
 & layered sweep & 6.30 s & 0.523 s & 51 & 54 & 53 & 0 & 0.523 s & 0.055 s & 0.012 s & 31.50 s & 4.5 GB \\
 & L-sweeps pad 8 & 3.62 s & 0.312 s & 53 & 56 & 55 & 0 & 0.312 s & 0.055 s & 0.012 s & 17.97 s & 4.3 GB \\
 & L-sweeps pad 16 & 5.22 s & 0.366 s & 53 & 56 & 55 & 0 & 0.366 s & 0.055 s & 0.012 s & 22.77 s & 5.4 GB \\
 & L-sweeps pad 8, inner 6 & 3.65 s & 1.975 s & 8 & 64 & 9 & 56 & 0.247 s & 0.055 s & 0.012 s & 16.21 s & 4.6 GB \\
256 & exact solve of $H$ & 31.91 s & 1.264 s & 10 & 11 & 10 & 0 & 1.264 s & 0.308 s & 0.046 s & 18.25 s & 16.6 GB \\
 & layered sweep & 31.40 s & 2.082 s & 156 & 164 & 163 & 0 & 2.082 s & 0.308 s & 0.046 s & 403.49 s & 16.4 GB \\
 & L-sweeps pad 8 & 12.46 s & 1.100 s & 328 & 345 & 344 & 0 & 1.100 s & 0.308 s & 0.046 s & 457.84 s & 14.6 GB \\
 & L-sweeps pad 16 & 16.97 s & 1.552 s & 317 & 333 & 332 & 0 & 1.552 s & 0.308 s & 0.046 s & 566.54 s & 18.4 GB \\
 & L-sweeps pad 24 & 23.72 s & 1.997 s & 297 & 312 & 311 & 0 & 1.997 s & 0.308 s & 0.046 s & 660.11 s & 23.3 GB \\
 & L-sweeps pad 8, inner 6 & 12.40 s & 7.673 s & 22 & 176 & 24 & 154 & 0.959 s & 0.308 s & 0.046 s & 177.83 s & 18.3 GB \\
\bottomrule
\end{tabular}\\[5pt]
\begin{minipage}{\textwidth}\footnotesize
\textbf{Columns.} $\omega/2\pi$: frequency; the grid is $n=8(\omega/2\pi)-1$ at eight
points per wavelength and the disc is $\omega/4\pi$ wavelengths across.
\emph{Inner solver}: the approximate inverse of $H$ the preconditioner uses; ``inner 6''
runs up to six GMRES iterations on $H$ per application, to $10^{-3}$, in place of a single
sweep. \emph{Setup}: building it --- sparsification, $H$, and the factorisation.
\emph{Apply}: one preconditioner application, timed once, comprising $Sr$, the inner solve
and the scatter. \emph{Outer}: iterations of the outer Krylov solver on the dense system.
\emph{Sweeps}: preconditioner applications performed inside that solve, LU solves on the
exact-$H$ row. \emph{Fwd}: forward maps, applications of $A=I+\omega^2K[m\,\cdot\,]$, each
an FFT convolution. \emph{Hmv}: applications --- not solves --- of the sparsified $H$, one
per inner GMRES iteration; only the two-level rows perform any.
$t_{\rm sw}$: the cost of one sweep, namely \emph{apply} divided by the sweeps that
application performs (one for the single-level rows, eight for the two-level rows from
$\omega/2\pi=16$ up); it carries $Sr$ and the scatter with it, $0.050$~s of $1.100$ at
$n=2047$. $t_{\rm fwd}$, $t_{\rm Hmv}$: the cost of one forward map and of one $H$ matvec.
Those two depend on the grid alone --- not on the medium, not on the preconditioner --- so
one measurement per size serves every row; $t_{\rm Hmv}$ is already inside the two-level
rows' $t_{\rm sw}$. \emph{Solve}: the outer Krylov solve, setup excluded. \emph{Peak}:
peak resident set of the process. Every column is measured.
\end{minipage}
\end{center}
The Mie error of the scattered field, relative $\ell^2$ on $20\,000$ grid points at eight
points per wavelength, is $0.16$, $0.10$, $0.074$, $0.076$, $0.083$, $0.20$, $0.45$ and
$0.56$ for $\omega/2\pi=2$ to $256$; at fixed $\omega/2\pi=8$ it falls as $0.074$,
$0.026$, $0.0090$, $0.0045$ at $8$, $16$, $32$, $64$ points per wavelength --- first order
in $h$, as a staircased discontinuous coefficient gives. So the solvers agree with each
other to their tolerance and the error is the discretisation's: at a fixed resolution it
grows with $\omega$ because the staircased interface shifts the effective radius by about
$h/2$ and a large, near-resonant scatterer is sensitive to that. The corrected
trapezoidal rule is built for smooth $m$; a sharp interface costs it about $64$ points
per wavelength for $5\%$ at $\omega/2\pi=256$.

The counts follow the solvers' schedule. A left-preconditioned GMRES(20) applies operator
and preconditioner once per iteration and once more at every restart, plus one
preconditioner application on the initial residual, so the single-level rows have
$\mathrm{fwd}=\mathrm{outer}+r$ and $\mathrm{sweeps}=\mathrm{outer}+1+r$ with
$r=\lfloor(\mathrm{outer}-1)/20\rfloor$ restarts, and the two-level rows have
$\mathrm{Hmv}=\mathrm{sweeps}-\mathrm{outer}$. Only the two-level sweep counts cannot be
derived: each application runs an inner GMRES that stops at $10^{-3}$ or its cap of six,
which is two sweeps per application at $\omega/2\pi\le8$ and eight from $16$ upward. As a
check, $\text{sweeps}\times t_{\rm sw}+\text{fwd}\times t_{\rm fwd}$ accounts for $392$~s
of the layered sweep's $403$ at $\omega/2\pi=256$, $176$ of the two-level scheme's $178$
and $17.0$ of the exact solve's $18.3$; it overshoots the single-level L-sweeps rows by
some $8\%$, so the single timed application those rows report is a little pessimistic
against the same sweep run back to back inside a solve. At $n=2047$ a forward map costs
$6.6$ times an $H$ matvec, $0.308$~s against $0.046$, and that ratio is the whole reason
the two-level scheme pays.

The iterations tell a different story from the lens. The exact solve of $H$ needs $3$,
$4$, $5$, $10$ iterations at $\omega/2\pi=32$ to $256$: the sparsification's floor is
modest. Everything above it is the sweeps' inexactness on a medium that reflects. The
lens reflects weakly and one pass, layered or L-shaped, captures almost the whole field;
the disc's interface sends waves back, each pass captures one more reflection, and the
layered sweep needs $156$ and the L-sweep $328$ --- with GMRES(20), eight and sixteen
restarts. It is not the absorber: a $3\lambda$ pad takes $328$ to $297$ while doubling the
application cost, where on the lens it took $14$ to $8$. Restarts cost too: $328$
iterations of GMRES(20) carry sixteen restarts, each an extra forward map and sweep. The
work columns show why the two-level scheme, useless on the lens, is decisive here. It does
not sweep less --- $176$ against the layered sweep's $164$ --- but its sweeps are L-sweeps
at about $0.96$~s rather than layered sweeps at $2.08$~s, and it cuts the forward maps from
$163$ to $24$; against single-level L-sweeps, $176$ sweeps against $345$ and $24$ forward
maps against $344$. The inner iterations move the work onto the sparse system $H$, where an
iteration costs a sweep and a sparse matvec, and off the dense system, where one costs a
forward map plus orthogonalisation against a basis of $4.2$-million-entry vectors: $190$~s
to solution against $435$ and $470$. The exact LU of $H$ is the limit of that trade and
beats all of them --- ten forward maps, eleven LU solves, $31.9$~s of setup against the
layered sweep's $31.4$, $50$~s in total at the same $16.6$~GB. In two dimensions at
these sizes the factorisation the sweeps exist to avoid is not expensive, and on a medium
that reflects, avoiding it costs an order of magnitude. The sweeps earn their keep where
that factorisation is prohibitive, which is three dimensions, and where the medium lets
one pass suffice.

\paragraph{Three dimensions.} The paper's own 3D construction is now in: a $q\times q$
decomposition of the $(x_1,x_2)$ plane, every box spanning the whole of $x_3$ with the
global PML at its ends --- their ``beam-shaped quasi-one-dimensional local problems'' ---
with the same eight directions and the same dataflow; trace layers become planes and the
diagonal trace an L-shaped prism. The scale differs: their Table~9 runs $q=12$ to $32$
beams per side on $p=q$ nodes, in constant media at six points per wavelength with
$2\lambda$ boxes and a $1\lambda$ PML; here $3\times3$ and $4\times4$ beams on one node,
$2.25\lambda$ boxes and a $0.5\lambda$ pad, in the lens. Converging lens, $h=\lambda/8$,
$b=4$, $C=12$ as in the package's 3D tables; L-sweeps on $16$ threads with compact
factors, the layered sweep untouched --- its application on one core, its 3D
factorisation using OpenBLAS's default $16$ threads in UMFPACK's dense fronts, which are
large enough here to matter ($7.5$ s of setup with one BLAS thread against $5.0$ s with
sixteen at $n=31$, unlike 2D where they only spun). ``In.$\times$out.'' is the nominal
count outer$\,\times\max(\mathrm{inner},1)$; ``sweeps'' is the number of sweep applications
actually performed inside the outer solve, and differs from it in both directions. A
left-preconditioned GMRES applies its preconditioner to the initial residual as well, so
the inner-free rows sweep outer$+1$ times; with inner iterations each of the outer
applications is an inner GMRES that stops at $10^{-3}$ after two iterations --- three
sweeps --- so those rows sweep about $3\times$outer times, half of inner$\times$outer:
\begin{center}\footnotesize\setlength{\tabcolsep}{4pt}
\begin{tabular}{rlrrcccrr}
\toprule
$\omega/2\pi$ & inner solver & setup & apply & outer & in.$\times$out. & sweeps & solve & peak \\
\midrule
4 & layered sweep                & 4.93 s   & 0.223 s & 4 & 4  & 5 & 1.20 s  & 2.5 GB \\
  & layered sweep, inner 6       & 4.97 s   & 0.716 s & 3 & 18 & 9 & 2.21 s  & 2.4 GB \\
  & L-sweeps $3\times3$, pad 4          & 10.54 s  & 0.161 s & 3 & 3  & 4 & 0.74 s  & 8.2 GB \\
  & L-sweeps $3\times3$, pad 4, inner 6 & 12.11 s  & 0.345 s & 3 & 18 & 8 & 1.33 s  & 8.1 GB \\
8 & layered sweep                & 28.26 s  & 1.748 s & 4 & 4  & 5 & 8.89 s  & 13.3 GB \\
  & layered sweep, inner 6       & 28.87 s  & 5.218 s & 3 & 18 & 9 & 15.97 s & 13.4 GB \\
  & L-sweeps $4\times4$, pad 4          & 169.90 s & 1.628 s & 4 & 4  & 5 & 8.36 s  & 39.9 GB \\
  & L-sweeps $4\times4$, pad 4, inner 6 & 162.59 s & 4.926 s & 3 & 18 & 9 & 15.01 s & 40.6 GB \\
\bottomrule
\end{tabular}
\end{center}
Inner iterations do not pay: the outer count sits at or next to the exact-$H$ floor of
three without them, so they only multiply the application cost. At $n=31$ L-sweeps is the
better inner solver --- three iterations against four, the application $1.4\times$ and the
solve phase $1.6\times$ faster --- at twice the setup and three times the memory. At
$n=63$ the layered sweep wins: the same four iterations, an application $7\%$ faster, the
solve phase within $6\%$, and a setup six times slower at three times the memory. The
reasons are geometric. A padded beam carries $(s+2(2+\mathrm{pad}))^2\,n_{\mathrm{tot}}$
unknowns for $s^2 n_{\mathrm{tot}}$ owned --- $3.1\times$ redundancy at $s=16$,
$\mathrm{pad}=4$, against $1.7\times$ for the 2D default --- 3D sparse LU cost grows like
$N^2$ rather than $N^{3/2}$, and pad 4 is already half a wavelength, so there is no thinner
absorber to reach for. The paper's estimate says the same --- 3D factorisation scales
like $n^3 n_{\mathrm{pml}}^6/p$, ``therefore it is crucial to keep the thickness of the
PML region as thin as possible'' --- and its own 3D applications take $52$~s at
$N=156^3$ on twelve nodes. A factorisation in flight holds about four times its final compact
size in UMFPACK transients; sixteen beams at once passed $45$ GB and were stopped, and
capping them at eight (\path{nfact = 8}) holds the peak to $39.9$ GB in two rounds.

\section{Outlook}

The paper demonstrates $O(N/p)$ scaling with one rank per row of subdomains, up to $128$
ranks at $12\,928$ points per side. On this one node the comparison in \S5 is not yet a
fair fight: L-sweeps wins on total time at every size measured, $1.96\times$ at
$\omega/2\pi=32$ falling to $1.35\times$ at $256$, but it wins by using sixteen cores
against a layered implementation that uses one. With the compact factors the margin at
$256$ is $1.69\times$ and the solve phase is ahead as well. In three dimensions the beam
variant loses to the layered sweep at $n=63$ on setup and memory, for the geometric
reasons given in \S5. And on a medium that reflects --- the constant-contrast disc ---
both sweeps need well over a hundred iterations at $256$ where the exact LU of $H$ needs
ten and costs no more to set up than the layered sweep: in two dimensions the case for
sweeping at all rests on the medium being smooth. Four things remain. The outer product should be distributed in the
same horizontal bands as the boxes, by a slab FFT or by a fast multipole method, after
which the all-reduce of the result disappears and only the Krylov inner products are
global. The row partition caps the useful rank count at the
number of box rows; a two-dimensional block partition is accepted by the code but its
communication pattern has not been tuned. The layered sweep's factorisation should be threaded before any of this is measured
again: its slice loop is independent per slice, and at the $3.1\times$ L-sweeps manages
here it would fall from $24.2$ s to about $7.3$ s at $\omega/2\pi=256$, enough to make the
layered sweep the faster method outright, since its solve phase already is. Until that is
done, \S5 compares implementations rather than algorithms. And the absorber wants a rule: held fixed in
wavelengths it lets the iteration count grow about linearly with the frequency at one
wavelength and roughly logarithmically at three. The box-size study rules out the
interface count as the driver, which leaves the size of the domain in wavelengths, and no
mechanism for that yet. A rule for the pad, and the reason behind it, is the first thing
worth measuring on a machine large enough to reach $\omega/2\pi=512$ and beyond.

\begin{thebibliography}{9}
\bibitem{LiuYing2018} F.~Liu, L.~Ying, Sparsify and sweep: an efficient preconditioner for
the Lippmann--Schwinger equation, SIAM J. Sci. Comput. 40 (2018) B379--B404.
\bibitem{Taus2020} M.~Taus, L.~Zepeda-N\'u\~nez, R.~J.~Hewett, L.~Demanet, L-Sweeps: a
scalable, parallel preconditioner for the high-frequency Helmholtz equation, J. Comput.
Phys. 420 (2020) 109706.
\bibitem{ZepedaZhao2016} L.~Zepeda-N\'u\~nez, H.~Zhao, Fast alternating bidirectional
preconditioner for the 2D high-frequency Lippmann--Schwinger equation, SIAM J. Sci.
Comput. 38 (2016) B866--B888.
\bibitem{Ying2015} L.~Ying, Sparsifying preconditioner for the Lippmann--Schwinger
equation, Multiscale Model. Simul. 13 (2015) 644--660.
\end{thebibliography}
\end{document}
